\clearpage\begingroup\SourceLeaf{402b}\ArithmeticHead{71}{R}{ELEMENTARE UND ANALYTISCHE GEOMETRIE.}
\renewcommand{\thefootnote}{\arabic{footnote})}\setcounter{footnote}{0}
\begin{center}
\textbf{23.}\par\smallskip
\textls[100]{Das Viereck.}
\end{center}
\noindent\hspace*{\parindent}Als \textsc{Gauss} die Zusätze zu \textsc{Schumachers} Übersetzung des \textsc{Carnot}schen
Werkes verfasste, löste er auch eine Aufgabe, die \textsc{Schumacher} im Oktober
1809 gestellt hatte, als \textsc{Gauss}, \textsc{Bessel} und er selbst ihren gemeinsamen Freund
\textsc{Olbers} in Bremen besuchten, die Aufgabe nämlich, in einem Viereck diejenige
Ellipse zu beschreiben, die den grössten möglichen Flächenraum umfasst.
\textsc{Schumacher} hatte sie den durch \textsc{Montucla} erneuerten \textit{Récréations mathématiques
et physiques} von \textsc{Ozanam} (Paris 1778) entnommen. Im Dezember 1809 wurde
\textsc{Gauss} von \textsc{Bessel} an die Aufgabe erinnert (Br. G.-\textsc{Bessel} S. 104). »Es ist
ein merkwürdiges Beispiel«, antwortet dieser am 7. Januar 1810, »wieviel bisweilen
von der Wahl der unbekannten Grössen abhängt. Ich setzte mich
gleich daran und kam, da ich zufällig hierin eine glückliche Wahl getroffen
hatte, sofort darauf, dass das ganze Problem bloss auf eine Gleichung zweiten
Grades sich reduziert« (Br. G.-\textsc{Bessel} S. 107). Die »glückliche Wahl« kam
darauf hinaus, dass er komplexe Grössen verwandte; in der Darstellung der
Lösung, die \textsc{Gauss} \textsc{Schumacher} mitteilte, ist dieser Ursprung zwar verhüllt
worden, aber doch noch deutlich genug sichtbar geblieben.
\par
Nachdem \textsc{Gauss} am 10. Februar 1810 an \textsc{Schumacher} geschrieben hatte,
\clearpage\SourceLeaf{403}\ArithmeticHead{72}{L}{STÄCKEL, GAUSS ALS GEOMETER.}
\noindent er habe eine sehr artige Auflösung gefunden und sei nicht abgeneigt, sie bekannt
zu machen (Br. G-\textsc{Sch.} I, S. 26), wurde die Aufgabe, wohl auf \textsc{Schumachers}
Veranlassung, im Maiheft der Monatlichen Correspondenz\footnote{Monatliche Correspondenz zur Beförderung der Erd- und Himmelskunde, herausgegeben von v. \textsc{Zach}, Bd. 21, 1810, S. 462.} den Mathematikern
vorgelegt, und das Augustheft brachte (S. 112—121) die Lösung von
\textsc{Gauss} (W. IV, S. 385). Im Besonderen wird darin der Lehrsatz bewiesen,
dass der geometrische Ort der Mittelpunkte der Ellipsen, die die vier Seiten
des Vierecks berühren, eine Gerade ist; daraus folgt als Zusatz, dass die
Mitten der drei Diagonalen eines Vierseits auf einer Geraden liegen.
\par
Das Septemberheft der Correspondenz enthält zwei weitere Lösungen, die
von \textsc{J. Fr. Pfaff} und \textsc{Mollweide} herrühren\footnote{Monatliche Correspondenz, Bd. 22, 1810, S. 223 und 227.}; eine vierte, von \textsc{Buzengeiger}
eingesandte konnte wegen Mangel an Raum nicht abgedruckt werden\footnote{A. a. O., S. 513.}. \textsc{Pfaff}
bemerkt, dass jener geometrische Ort schon bei \textsc{Newton}\footnote{I. \textsc{Newton}, \textit{Philosophiae naturalis principia mathematica}, London 1687, Liber I, Lemma 25; im Corollarium 3 wird auch der Satz ausgesprochen, dass die Mitten der Diagonalen eines Vierseits auf einer Geraden liegen.} und \textsc{Euler}\footnote{L. \textsc{Euler}, \textit{Introductio in analysin}, T. II, Lausanne 1748, § 123.} zu
finden sei. Endlich gab \textsc{Schumacher} im Novemberheft\footnote{Monatliche Correspondenz, Bd. 22, 1810, S. 505.} eine Ergänzung,
indem er zeigte, dass unter Umständen eine innerhalb des Vierecks liegende
Ellipse, die nur drei Seiten berührt, den grössten Inhalt liefert. Die Aufgabe
ist später wiederholt bearbeitet worden; \textsc{Plücker}, \textsc{Schläfli} und \textsc{Steiner} haben
sich um sie bemüht\footnote{J. \textsc{Plücker}, \textit{Analytisch-geometrische Entwicklungen}, Band II, Essen 1831, S. 208; L. \textsc{Schläfli}, \textit{Anwendungen des barycentrischen Calculs}, Archiv der Mathematik und Physik, Bd. 12, 1849, S. 99; J. \textsc{Steiner}, \textit{Teoremi relativi alle coniche inscritte e circoscritte}, Giornale arcadico, t. 99, S. 147, \textsc{Crelles} Journal, Bd. 30, 1845, S. 17, Gesammelte Werke, Bd. II, S. 334. \textsc{Euler} hat die duale Aufgabe behandelt, um ein gegebenes Viereck die kleinste Ellipse zu beschreiben, Nova acta acad. sc. Petrop. 9 (1791), 1795, S. 132; vorgelegt den 4. Sept. 1777.}.
\par
In die Zeit um 1810 gehört auch wohl eine Aufzeichnung, die sich auf
der letzten Seite des \textsc{Gauss}schen Exemplares des ersten Teiles der \textsc{Schumacher}schen
Übersetzung befindet. \textsc{Carnot} hatte in einer 1806 erschienenen Abhandlung,
die \textsc{Schumacher} in seine Ausgabe aufgenommen hat, die zwischen
den Seiten und den Diagonalen eines Vierecks bestehende Gleichung
\clearpage\SourceLeaf{404}\ArithmeticHead{73}{R}{ELEMENTARE UND ANALYTISCHE GEOMETRIE.}\setcounter{footnote}{0}
\noindent hergeleitet (2. Teil, S. 258). und \textsc{Gauss} gibt einen einfachen Beweis dieser für die
Ausgleichungsrechnungen der Geodäsie wichtigen Beziehung (W. IX, S. 248).
\par
Mit den geodätischen Messungen, die \textsc{Gauss} von 1821 bis 1825 anstellte,
hängt es auch zusammen, dass er sich eingehend mit einer Aufgabe beschäftigt
hat, die nach einem Mathematiker, der weder ihr Urheber noch ihr erster
Löser ist, der sich aber Verdienste um sie erworben hat, häufig als \textsc{Pothenot}sches
Problem bezeichnet wird\footnote{L. \textsc{Pothenot}, \textit{Problème de Géométrie pratique}, Mém. de l’Acad. depuis 1666 jusqu’à 1699, t. 10, Paris 1730, S. 150 (vorgelegt 1692). Für das Geschichtliche vgl. die Dissertation von R. \textsc{Wagner}, \textit{Über das Pothenotsche Problem}, Göttingen 1852 und die Angaben in J. C. \textsc{Poggendorffs} Biographisch-literarischem Handwörterbuch, II, Leipzig 1863, Spalte 509.}. Es handelt sich darum, bei einer trigonometrischen
Aufnahme die Lage eines Punktes dadurch festzulegen, dass die
Winkel gemessen werden, welche die von ihm nach drei bekannten Punkten
(Netzpunkten) gehenden Richtungen miteinander bilden (Rückwärtseinschneiden).
Die im Nachlass von \textsc{Gauss} befindlichen umfangreichen Aufzeichnungen über
das \textsc{Pothenot}sche Problem aus den Jahren 1832 bis 1852 werden ergänzt
durch Briefe an \textsc{Gerling} und \textsc{Schumacher} aus den Jahren 1830 bis 1842 und
durch die Ausarbeitung einer im Jahre 1840 gehaltenen Vorlesung \textit{über die
Theorie der imaginären Grössen} (W. VIII, S. 307—334).
\par
Die Heranziehung der komplexen Grössen erweist sich hier als besonders
nützlich. Indem \textsc{Gauss} den Ecken \(a_0,a_1,a_2,a_3\) des Vierecks, das aus dem
festzulegenden Punkt und den drei Netzpunkten besteht, das Dreieck zuordnet,
dessen Ecken durch die aus der Lehre von den biquadratischen Gleichungen
wohlbekannten Verbindungen
\[
a_0a_1+a_2a_3,\quad a_0a_2+a_3a_1,\quad a_0a_3+a_1a_2
\]
bestimmt werden, gelangt er zu seiner »zierlichen Auflösung«; zu demselben
Dreieck war übrigens schon \textsc{Collins} durch einen geometrischen Kunstgriff
gelangt\footnote{J. \textsc{Collins}, \textit{A solution of a chorographical problem}, Philosophical transactions, Vol. 6, Nr. 69, London, März 1671.}.
\par
\textsc{Gauss} eigentümlich ist die Frage nach der »physischen Möglichkeit der
Daten in \textsc{Pothenots} Aufgabe«. Wenn nämlich beim Rückwärtseinschneiden
der Punkt gesucht wird, von dem aus zwei gegebene, aneinander stossende
Strecken unter gemessenen Winkeln erscheinen, so ist man sicher, dass die
\par\smallskip\noindent{\scriptsize X\textsubscript{2} Abh. 4.\hfill 10}\par
\clearpage\SourceLeaf{405a}\ArithmeticHead{74}{L}{STÄCKEL, GAUSS ALS GEOMETER.}\setcounter{footnote}{0}
\noindent Aufgabe, sobald nur der gefährliche Kreis vermieden wird, eine bestimmte
Lösung hat. Anders steht es, wenn jene Winkel willkürlich angenommen
werden. Dann braucht es keine Lösung zu geben, und es entsteht die Frage
nach einem Kennzeichen für die Lösbarkeit; \textsc{Gauss} hat darauf eine überraschend
einfache Antwort gegeben.
\par
Warum, wird man fragen, hat \textsc{Gauss} auf einen so elementaren Gegenstand
so viel Zeit und Mühe verwendet? Aufschluss hierüber gibt der Brief
an \textsc{Gerling} vom 14. Januar 1842. Nachdem er diesem das Kennzeichen mitgeteilt
hat, bittet er ihn, es für sich zu behalten, »weil ich das Theorem,
womit es zusammenhängt, selbst einmal bei schicklicher Gelegenheit zu behandeln
mir vorbehalte, weniger wegen der Eleganz des Theorems an sich,
als wegen der Eleganz, welche die Anwendung der komplexen Grössen dabei
darbietet, also namentlich bei einer Gelegenheit, wo ich mehr von dem Gebrauch
der komplexen Grössen sagen kann« (W. VIII, S. 315)\footnote{Für die Behandlung geometrischer Aufgaben mittels komplexer Grössen vgl. noch die Dissertation von H. \textsc{zur Nedden}, \textit{Applicatio numeri complexi ad demonstranda nonnulla geometriae theoremata}, Göttingen 1840.}.
\clearpage\endgroup
